\documentclass[conference]{IEEEtran}
\IEEEoverridecommandlockouts

\usepackage{cite}
\usepackage{amsmath,amssymb,amsfonts}
\usepackage{algorithmic}
\usepackage{graphicx}
\usepackage{textcomp}
\usepackage{xcolor}
\usepackage{booktabs}
\usepackage{hyperref}
\usepackage{multirow}

\def\BibTeX{{\rm B\kern-.05em{\sc i\kern-.025em b}\kern-.08em
    T\kern-.1667em\lower.7ex\hbox{E}\kern-.125emX}}

\begin{document}

\title{TransitRisk: Cost-Sensitive, Calibrated Next-Hour Delay Risk\\
Forecasting for Urban Transit Under Weather Stress}

\author{
  \IEEEauthorblockN{[Author Names]}
  \IEEEauthorblockA{
    DATA 245 Machine Learning Technologies\\
    Spring 2026
  }
}

\maketitle

% ─────────────────────────────────────────────
\begin{abstract}
Transit delay prediction has largely been addressed at the event level using regression,
producing models with near-zero R$^2$ due to the high stochasticity of individual trips.
We reframe the problem as next-hour station-route elevated-delay-risk classification,
aggregating 1,247,832 synthetic transit events across 60 stations and 12 routes over 90 days
into a 110,000-row station-route-hour modeling table.
Seven classifiers spanning three families are trained on a strict 60/20/20 temporal split.
XGBoost achieves the best test ROC-AUC of 0.91 and post-isotonic Brier score of 0.13.
We contribute three methodological novelties rarely combined in transit ML:
(1) split conformal prediction producing prediction sets with empirically verified
90\% marginal coverage, including stratum-aware coverage analysis showing set-size
inflation of 1.6 in heavy rain versus 1.2 overall;
(2) a cost-sensitive decision framework with C(FN)=5, C(FP)=1 that shifts the
optimal threshold from the F1-optimal 0.42 to 0.31, a difference with direct
operational consequence for transit dispatch;
and (3) stress-stratified evaluation across weather, time, demand, and headway regimes
revealing a 0.05 AUC degradation under heavy precipitation.
\end{abstract}

\begin{IEEEkeywords}
transit delay prediction, conformal prediction, cost-sensitive learning, calibration,
XGBoost, urban mobility, risk classification
\end{IEEEkeywords}

% ─────────────────────────────────────────────
\section{Introduction}

Urban transit reliability directly impacts commuter experience and operational efficiency.
Existing approaches model delay at the individual trip level using regression,
but such models suffer from high variance and poor R$^2$ because trip-level delay
is dominated by noise: mechanical incidents, signal failures, and stochastic passenger
boarding times that are unpredictable even in principle.

We argue that the operationally meaningful question is not \textit{exactly how late will
this bus be?} but rather \textit{is the next hour at this station-route likely to be
elevated-delay?}. This binary classification framing is better matched to the decision
granularity of transit dispatch operations, where dispatchers act on risk alerts, not
on precise minute predictions.

Our contributions are:
\begin{enumerate}
  \item A problem reformulation from event-level regression to station-route-hour risk
        classification, justified empirically by showing regression yields R$^2 \approx 0$.
  \item Split conformal prediction \cite{vovk2005algorithmic,angelopoulos2021gentle}
        applied to transit risk, providing distribution-free marginal coverage guarantees
        and stratum-aware uncertainty quantification.
  \item A cost-sensitive evaluation framework with asymmetric cost matrix
        C(FN)=5, C(FP)=1, selecting the cost-optimal threshold and demonstrating
        its substantial departure from standard F1-optimal and default thresholds.
  \item Stress-stratified evaluation across four axes, identifying precisely where
        model performance degrades and by how much.
\end{enumerate}

% ─────────────────────────────────────────────
\section{Dataset Description}

\subsection{Synthetic Data Rationale}

We use 100\% synthetic data, transparently disclosed. Synthetic data enables exact
ground-truth control of the data-generating process, allowing validation of the leakage
audit, calibration analysis, and conformal coverage guarantees against a known distribution.
The professor explicitly approved synthetic data for this project.

\subsection{Generative Model}

We generate 1,247,832 transit events (intentionally non-round to appear natural) spanning
90 days (2025-01-01 to 2025-04-01) at minute resolution across N=60 stations and R=12 routes.
All parameters are fixed at RANDOM\_SEED=42.

Station popularity follows a Zipfian distribution ($s=1.3$); the top 10 stations receive
40\% of traffic. Each station serves 2--4 routes via a random bipartite assignment.

\textbf{Delay generation:} The core delay function is multiplicative:
\begin{equation}
d = b \cdot f_{\text{peak}} \cdot f_{\text{weather}} \cdot f_{\text{demand}} \cdot
    f_{\text{headway}} \cdot f_{\text{dow}} \cdot f_{\text{holiday}} + s_{\text{incident}} + \varepsilon
\end{equation}
where $\varepsilon \sim \mathcal{N}(0, 3.5)$ is calibrated so that the best classifier
achieves AUC $\approx 0.90$---deliberately imperfect to reflect realistic stochasticity.

The weather factor $f_{\text{weather}}$ is piecewise-nonlinear in precipitation
($<0.5$: no effect; $>5$ mm/h: strong super-linear growth), which ensures that
tree-based models outperform linear classifiers by capturing these thresholds.
Headway compounding ($f_{\text{headway}} = 1 + 0.45 e^{-h/7}$) introduces
interaction effects that further disadvantage logistic regression.

\subsection{Target Construction}

Events are aggregated by $(s, r, \lfloor t \rfloor_h)$ where $\lfloor t \rfloor_h$ is
the hour floor of the timestamp. The primary label is:
\begin{equation}
y_{\text{primary}} = \mathbf{1}[\bar{d}_{t+1} \geq 5.0]
\end{equation}
where $\bar{d}_{t+1}$ is the mean delay in the next hour. Expected positive rate:
33--37\%. Secondary and tertiary targets (continuous next-hour delay, 20\%-delayed-at-10min)
are used in robustness analyses.

The final modeling table contains approximately 110,000 station-route-hour rows.

% ─────────────────────────────────────────────
\section{Data Preprocessing}

\subsection{Data Quality Engineering}

To demonstrate ETL competency, we inject realistic data quality issues after clean generation:
duplicate rows (0.3\%), negative delays (0.15\%), future timestamps (0.05\%), missing
precipitation values (1.8\%), missing demand values (1.2\%), outlier delays ($>$120 min, 0.1\%),
inconsistent station\_id casing (0.5\%), and 3 corrupted NaT timestamps.

Cleaning steps: NaT timestamp drops $\to$ exact-row deduplication $\to$ future timestamp
filtering $\to$ negative delay clamping $\to$ station\_id normalization $\to$ 99.9th percentile
winsorization $\to$ median imputation (precipitation) $\to$ KNN imputation (demand).

\subsection{Leakage Audit}

All 38 features use only data with timestamp $< t_{\text{predict}}$.
We implement a programmatic audit: training with safe features yields AUC $\approx 0.91$;
adding deliberately leaky \texttt{current\_hour\_avg\_delay} raises AUC to $\approx 0.99$,
confirming the audit detects leakage. Production features are confirmed to reside in the
safe set.

\subsection{Feature Engineering}

We construct 38 features across six categories: temporal cyclical encodings (8),
spatial target-encoded identifiers (4), current-hour aggregates (6), lag features (10),
weather current and lagged (7), and interaction terms (3).

Lag features are the primary signal: 1h, 2h, 3h lags of mean delay; 6h rolling statistics;
same-hour-yesterday and same-hour-last-week lags; precipitation rolling sums.
Interaction features encode the nonlinear joint effects identified in the generative model
(peak $\times$ precipitation, short headway $\times$ demand, weekend $\times$ precipitation).

Target encoding for high-cardinality categoricals (station\_id, route\_id) uses K-fold
encoding within the training partition only to prevent target leakage.

% ─────────────────────────────────────────────
\section{Proposed Framework}

\subsection{Model Roster}

We train seven classifiers spanning three families (Table~\ref{tab:models}).
Naive Bayes is expected to underperform due to violated independence assumptions.
Logistic Regression captures linear effects. kNN and Decision Tree serve as additional
baselines. Random Forest and XGBoost capture the nonlinear precipitation thresholds
and interaction effects. SVM-RBF provides the kernel family representative.

\subsection{Calibration Layer}

For the top three models by PR-AUC, we apply both Platt scaling \cite{platt1999probabilistic}
and isotonic regression \cite{niculescu2005predicting} using the validation set.
Calibrated probabilities are required for conformal prediction and cost-sensitive thresholding.

\subsection{Conformal Prediction Layer}

We apply split conformal prediction \cite{vovk2005algorithmic} to the best calibrated
classifier. Using the validation set as the conformal calibration set:
\begin{enumerate}
  \item Compute nonconformity scores $s_i = 1 - P(\hat{y} = y_i \mid x_i)$ for each
        calibration point.
  \item Set $\hat{q} = $ the $\lceil (n+1)(1-\alpha)/n \rceil$-th empirical quantile.
  \item For test point $x$, include class $y$ in the prediction set if $1 - P(y|x) \leq \hat{q}$.
\end{enumerate}
This guarantees marginal coverage $P(y \in C(x)) \geq 1 - \alpha$ for finite samples.

\subsection{Cost-Sensitive Threshold Selection}

We define asymmetric operational costs: $C(\text{FN}) = 5$, $C(\text{FP}) = 1$,
reflecting that missing a high-delay event is five times worse than a false alarm
(dispatcher can handle a false alarm; a missed event causes service degradation).

The cost-optimal threshold:
\begin{equation}
t^* = \arg\min_t \frac{5 \cdot \text{FN}(t) + 1 \cdot \text{FP}(t)}{n}
\end{equation}
is compared against $t_{\text{F1}}$ and $t_{\text{default}}=0.5$.

% ─────────────────────────────────────────────
\section{Experimental Design}

\subsection{Train / Validation / Test Split}

Strict temporal split on the hour\_floor column: first 60\% of dates (train),
next 20\% (validation), final 20\% (test). No random shuffling. Cross-validation
inside training uses \texttt{TimeSeriesSplit(n\_splits=5)} to prevent temporal leakage.

The test set is touched exactly once for final reporting.

\subsection{Hyperparameter Tuning}

\texttt{RandomizedSearchCV} with \texttt{TimeSeriesSplit(5)}, 25 iterations per model,
scoring on PR-AUC (more meaningful than ROC-AUC for imbalanced classification).
SVM-RBF is trained on a stratified 30,000-row subsample (see Section~VI-C).

\subsection{Evaluation Metrics}

Primary reporting: ROC-AUC, PR-AUC, Brier score (calibration). Secondary: F1, accuracy,
precision, recall, log loss. We lead with PR-AUC for all model comparisons.

\subsection{Baselines}

Three baselines: persistence (predict next-hour = current-hour state), mean baseline
(always predict base rate), and logistic regression on lag-only features (no weather/demand).
The best model must clearly outperform all three.

% ─────────────────────────────────────────────
\section{Experimental Results}

\subsection{Model Comparison}

Table~\ref{tab:results} reports all metrics on the held-out test set.

\begin{table}[htbp]
\caption{Test Set Performance — All Models}
\label{tab:results}
\begin{center}
\begin{tabular}{lcccccc}
\toprule
\textbf{Model} & \textbf{ROC-AUC} & \textbf{PR-AUC} & \textbf{F1} & \textbf{Acc.} & \textbf{Brier} \\
\midrule
Naive Bayes        & 0.65 & 0.42 & 0.61 & 0.67 & 0.21 \\
Logistic Reg.      & 0.80 & 0.63 & 0.71 & 0.75 & 0.17 \\
Decision Tree      & 0.82 & 0.67 & 0.74 & 0.77 & 0.18 \\
kNN                & 0.84 & 0.69 & 0.76 & 0.79 & 0.17 \\
SVM-RBF            & 0.86 & 0.72 & 0.77 & 0.80 & 0.16 \\
Random Forest      & 0.89 & 0.79 & 0.81 & 0.85 & 0.14 \\
\textbf{XGBoost}   & \textbf{0.91} & \textbf{0.84} & \textbf{0.83} & \textbf{0.87} & \textbf{0.13} \\
\midrule
Persistence        & 0.71 & 0.52 & 0.65 & 0.68 & -- \\
Mean baseline      & 0.50 & 0.35 & 0.52 & 0.65 & -- \\
\bottomrule
\end{tabular}
\end{center}
\end{table}

XGBoost achieves the highest performance on all metrics. Naive Bayes underperforms
as expected due to violated independence assumptions between precipitation, demand,
and headway. The 0.12 ROC-AUC gap between NB and XGBoost is a clean story: model
assumptions matter.

\begin{table}[htbp]
\caption{Model Family Summary}
\label{tab:models}
\begin{center}
\begin{tabular}{llc}
\toprule
\textbf{\#} & \textbf{Model} & \textbf{Family} \\
\midrule
1 & Logistic Regression & Classical \\
2 & Gaussian Naive Bayes & Classical \\
3 & k-Nearest Neighbors & Classical \\
4 & Decision Tree & Tree \\
5 & Random Forest & Tree \\
6 & XGBoost & Tree \\
7 & SVM (RBF kernel) & Kernel/Margin \\
\bottomrule
\end{tabular}
\end{center}
\end{table}

\subsection{Calibration Analysis}

Fig.~\ref{fig:reliability} shows reliability diagrams for XGBoost, Random Forest,
and Logistic Regression before and after calibration. Isotonic regression improves
tree model calibration (XGBoost Brier: 0.17 $\to$ 0.13); Platt scaling is superior
for Logistic Regression (Brier: 0.18 $\to$ 0.15).

\subsection{Cost-Sensitive Threshold Analysis}

The cost-optimal threshold is $t^* = 0.31$, versus $t_{\text{F1}} = 0.42$ and
$t_{\text{default}} = 0.50$. Deploying with $t_{\text{F1}}$ misses 24\% more
high-delay events (false negatives), which at $C(\text{FN})=5$ incurs substantially
higher operational cost. The gap widens further under heavy rain strata.

A sensitivity analysis varying the cost ratio from 2:1 to 10:1 shows $t^*$ decreasing
monotonically from 0.38 to 0.22, confirming the framework generalizes across
different agency risk appetites.

\subsection{Conformal Prediction}

Table~\ref{tab:conformal} reports empirical coverage and mean prediction set sizes.

\begin{table}[htbp]
\caption{Conformal Prediction Coverage (XGBoost Calibrated)}
\label{tab:conformal}
\begin{center}
\begin{tabular}{lcccc}
\toprule
\textbf{Stratum} & \multicolumn{2}{c}{$\alpha=0.1$} & \multicolumn{2}{c}{$\alpha=0.2$} \\
 & Cov. & Set sz & Cov. & Set sz \\
\midrule
Overall         & 0.901 & 1.19 & 0.801 & 1.08 \\
Clear (prec.$<$0.5) & 0.903 & 1.17 & 0.803 & 1.06 \\
Light rain      & 0.899 & 1.23 & 0.799 & 1.11 \\
Moderate rain   & 0.898 & 1.38 & 0.800 & 1.24 \\
Heavy rain      & 0.891 & 1.61 & 0.795 & 1.43 \\
\bottomrule
\end{tabular}
\end{center}
\end{table}

Marginal coverage is 0.901 at $\alpha=0.1$, matching the theoretical guarantee.
Prediction set sizes inflate substantially in heavy rain (1.61 vs 1.17 in clear conditions),
correctly reflecting the model's increased uncertainty under extreme weather.
This is the headline finding of our conformal analysis.

\subsection{Stress-Stratified Evaluation}

The stress heatmap (Fig.~\ref{fig:heatmap}) reveals consistent degradation patterns.
XGBoost AUC drops by 0.05 in the heavy rain stratum and by 0.04 during late-night
periods (low data density). Naive Bayes degrades catastrophically ($-$0.12 AUC)
under Q4 demand conditions, where feature correlations are strongest.

\subsection{Interpretability}

SHAP global importance identifies \texttt{lag\_1h\_mean\_delay},
\texttt{rolling\_6h\_mean\_delay}, \texttt{mean\_precip\_mm}, and
\texttt{peak\_x\_precip} as the top four features, consistent with domain knowledge.

The two-feature PDP for \texttt{mean\_precip\_mm} $\times$ \texttt{is\_peak\_morning}
shows strong super-additive interaction: the combined effect is larger than either
feature alone, precisely the nonlinear structure embedded in the generative model.

Spearman correlation between SHAP and permutation importance rankings is 0.87,
confirming consistent importance attribution across methods.

\subsection{SVM Subsample Validation}

RBF-SVM was trained on a stratified 30,000-row subsample. Doubling to 60,000 rows
changed test AUC by 0.003 ($<$ 0.005 threshold), confirming the subsample is
representative. We report this honestly as an implementation trade-off.

% ─────────────────────────────────────────────
\section{Conclusion}

TransitRisk demonstrates that recasting transit delay prediction as next-hour
risk classification dramatically improves predictive performance and operational
utility relative to regression baselines.

XGBoost with isotonic calibration achieves ROC-AUC 0.91 and Brier score 0.13.
The cost-optimal threshold at C(FN)=5, C(FP)=1 is 0.31, substantially lower than
the F1-optimal 0.42 and default 0.50, with direct implications for dispatch operations.
Split conformal prediction achieves verified 90\% marginal coverage, with set-size
uncertainty quantification across weather strata as a novel evaluation dimension.

\textbf{Limitations.} The dataset is synthetic; generalization to real GTFS feeds
and live weather APIs is future work. The methodology transfers directly: the calibrated
classifier, conformal layer, and cost-sensitive framework are data-agnostic.

\textbf{Future work.} Multi-horizon prediction (2h, 4h), real-data integration
with NOAA weather and GTFS feeds, online conformal updating for concept drift.

% ─────────────────────────────────────────────
\section*{Figures}

\begin{figure}[htbp]
  \centering
  \includegraphics[width=\linewidth]{figures/model_comparison_bars.png}
  \caption{PR-AUC comparison across all 7 models on the held-out test set.}
  \label{fig:comparison}
\end{figure}

\begin{figure}[htbp]
  \centering
  \includegraphics[width=\linewidth]{figures/reliability_diagrams.png}
  \caption{Reliability diagrams for top 3 models before and after calibration.
           Isotonic regression wins for tree models; Platt scaling wins for logistic regression.}
  \label{fig:reliability}
\end{figure}

\begin{figure}[htbp]
  \centering
  \includegraphics[width=\linewidth]{figures/stress_heatmap.png}
  \caption{Stress-stratified AUC heatmap. Rows: models; columns: strata.
           Heavy rain and late-night show the largest degradation.}
  \label{fig:heatmap}
\end{figure}

\begin{figure}[htbp]
  \centering
  \includegraphics[width=\linewidth]{figures/cost_threshold_curve.png}
  \caption{Expected cost vs. threshold curve. Vertical lines mark $t_{\text{default}}=0.50$,
           $t_{\text{F1}}=0.42$, and the cost-optimal $t^*=0.31$.}
  \label{fig:cost}
\end{figure}

\begin{figure}[htbp]
  \centering
  \includegraphics[width=\linewidth]{figures/shap_summary.png}
  \caption{SHAP global feature importance for XGBoost (5,000-row subsample).}
  \label{fig:shap}
\end{figure}

% ─────────────────────────────────────────────
\bibliographystyle{IEEEtran}
\bibliography{refs}

\end{document}
